\subsection{The Eberlein and Šmulian theorems}

THEOREM 1 (Eberlein). --- Let E be a Hausdorff and quasi-complete locally convex space, T a topology on E which is compatible with the duality between E and E' and A a subset of E. For A to be relatively compact for T, it is necessary and sufficient that every infinite sequence of points of A has a limit point in E for T.

The condition stated is obviously necessary.

Suppose now that every infinite sequence of points of A has a limit point for T, hence also for the coarser topology \(\sigma(\mathrm{E}, \mathrm{E}')\) . Then A is precompact for T (IV, p.~32, prop. 1); in order that A be relatively compact for T, it is necessary and sufficient that it be so for \(\sigma(\mathrm{E}, \mathrm{E}')\) (loc. cit.). Therefore it is enough to prove the theorem when T is the weakened topology \(\sigma(\mathrm{E}, \mathrm{E}')\) .

Let \(\hat{\mathrm{E}}\) denote the completion of E, which we shall identify as usual with a subspace of the algebraic dual \(\mathrm{E}^{\prime *}\) of \(\mathrm{E}'\) (III, p.~21, th. 2). Let \(\mathrm{E}_{\sigma}\) , \(\hat{\mathrm{E}}_{\sigma}\) and \(\mathrm{E}_{\sigma}^{\prime *}\) denote the spaces E, \(\hat{\mathrm{E}}\) and \(\mathrm{E}^{\prime *}\) endowed with the topologies \(\sigma (\mathrm{E},\mathrm{E}')\) , \(\sigma (\hat{\mathrm{E}},\mathrm{E}')\) and \(\sigma (\mathrm{E}^{\prime *},\mathrm{E}^{\prime})\) respectively.

Let \((x_{i}^{\prime})_{i\in\mathbb{I}}\) be a basis of the vector space \(E^{\prime}\) over the field K. The mapping \(f\mapsto(f(x_{i}^{\prime}))_{i\in\mathbb{I}}\) is a homeomorphism \(\phi\) from \(E_{\sigma}^{\prime*}\) onto \(K^{I}\) ; for every \(i\in I\) , the image of A under the mapping \(x_{i}^{\prime}\) from E into K is relatively compact: for, K is metrizable and every infinite sequence of elements of \(x_{i}^{\prime}(A)\) has a limit point. It follows that \(\phi(A)\) is relatively compact in \(K^{I}\) , hence that the closure \(\overline{A}\) of A in \(E_{\sigma}^{\prime*}\) is compact.

Next we shall prove that \(\overline{A}\) is contained in \(\hat{E}\) . Let H be an equicontinuous subset of \(E'\) ; let X be its closure for \(\sigma(E', E)\) ; X is compact (III, p.~17, cor. 2). For every \(x \in E'^{*}\) , let \(\phi_{x}\) be the restriction of \(x' \mapsto \langle x, x' \rangle\) to X; let \(\tilde{A} \subset \mathcal{C}_{s}(X)\) be the set of functions \(\phi_{x}\) as x ranges over A. In view of the hypothesis on A, every infinite sequence of elements of \(\tilde{A}\) has a limit point in \(\mathcal{C}_{s}(X)\) ; by prop. 2 (IV, p.~33), the set \(\tilde{A}\) is therefore relatively compact in \(\mathcal{C}_{s}(X)\) . It follows that for every \(a \in \overline{A}\) , the function \(\phi_{a}\) on X is continuous. The inclusion \(\overline{A} \subset \hat{E}\) then follows from th. 2 of III, p.~21.

Now we shall show that \(\overline{\mathbf{A}}\) is contained in E. Since A is precompact in \(\mathrm{E}_{\sigma}\) (IV, p.~32, prop. 1), it is bounded in \(\mathrm{E}_{\sigma}\) (III, p.~3, prop. 2), hence also in E (IV, p.~1, prop. 1). Let C be the closed convex balanced envelope of A in E. Then C is bounded since A is bounded, hence complete since E is quasi-complete. In other words, C is a convex and closed subset of \(\hat{\mathrm{E}}\) , so also of \(\hat{\mathrm{E}}_{\sigma}\) (IV, p.~1, prop. 1). Since \(\mathbf{A} \subset \mathbf{C}\) and the topology of \(\hat{\mathrm{E}}_{\sigma}\) is induced by that of \(\mathrm{E}_{\sigma}^{\prime *}\) , we have \(\overline{\mathbf{A}} \subset \mathbf{C}\) , and hence \(\overline{\mathbf{A}} \subset \mathbf{E}\) .

Since the topology of \(E_{\sigma}\) is induced by that of \(E_{\sigma}^{\prime*}\) , the subset \(\overline{A}\) of \(E_{\sigma}\) is compact, and th. 1 follows.

THEOREM 2 (Šmulian). --- Let E be a Fréchet space and A a subset of E. Let \(E_{\sigma}\) denote the space E endowed with the weakened topology. The following conditions are equivalent :

\begin{enumerate}
\def\labelenumi{(\roman{enumi})}
\item
  A is relatively compact in \(\mathrm{E}_{\sigma}\) ;
\item
  every infinite sequence of points of A has a limit point in \(\mathrm{E}_{\sigma}\) ;
\item
  from every infinite sequence of points of \(\mathbf{A}\) , we can extract a sequence which converges in \(\mathrm{E}_{\sigma}\) .
\end{enumerate}

The equivalence of (i) and (ii) follows from Eberlein's theorem and (iii) obviously implies (ii).

We shall prove that (i) implies (iii). Suppose that the closure B of A in \(E_{\sigma}\) is compact and that \((x_{n})_{n\in\mathbf{N}}\) is a sequence of points of A. Let F denote the smallest closed vector subspace of E containing the \(x_{n}\) , this is a Fréchet space satisfying the first axiom of countability. Since F is closed in \(E_{\sigma}\) and the topology \(\sigma(\mathrm{F},\mathrm{F}')\) on F is induced by \(\sigma(\mathrm{E},\mathrm{E}')\) , the set \(B\cap F\) is compact for \(\sigma(\mathrm{F},\mathrm{F}')\) . On account of the remarks in IV, p.~32, the existence of a sequence extracted from \((x_{n})_{n\in\mathbf{N}}\) converging for \(\sigma(\mathrm{E},\mathrm{E}')\) (or, which is the same, for \(\sigma(\mathrm{F},\mathrm{F}')\) ) is a consequence of the following lemma:

Lemma 1. --- Let F be a Fréchet space satisfying the first axiom of countability. Every subset C of F which is compact for the topology T induced by \(\sigma(\mathrm{F},\mathrm{F}')\) is metrizable for this topology.

Since the topology of precompact convergence on \(F'\) is finer than the topology \(\sigma(F', F)\) , there exists an everywhere dense countable subset in \(F_{s}'\) (III, p.~18, cor. 1). Hence the set C can be identified with a subset of \(K^{D}\) , and the topology induced on C by that of \(K^{D}\) , which is metrizable (GT, IX, § 2, No.~8) is coarser than the topology induced by \(\sigma(F, F')\) , for which C is compact. Hence these two topologies are identical (GT, I, § 9, No.~4, cor. 3).

Šmulian's theorem can be extended to the case where E is the strict inductive limit of a sequence of Fréchet spaces (IV, p.~67, exerc. 2).

